%%%%% PRELIMINARES %%%%%

%% Plantilla para la formulacion de presentaciones de estudiantes o profesores de la Facultad de Economia, de la Universidad Peruana de Ciencias Aplicadas. De no pertenecer a la comunidad academica, sirvase abstenerse.

%% FORMATO DEL DOCUMENTO
\documentclass[fleqn]{beamer}
\usepackage[spanish]{babel}
\selectlanguage{spanish}
\usepackage[utf8]{inputenc}
\usepackage{framed}
\usetheme{ZMBZFMK}
\usefonttheme[onlysmall]{structurebold}
\mode<presentation>
%\setbeamercovered{transparent=10}

%% HABILITA USO DE MATRICES EN ECUACIONES
\newcommand{\projdim}[1]{\mathbb{P}^#1}
\newcommand{\esperanza}[1]{\mathbb{E}(#1)}
\newcommand{\proba}[1]{\mathbb{P}(#1)}
\newcommand{\gavilla}[2]{\mathcal{O}_{#1}(#2)}

%% HABILITA USO DE DIAGRAMAS
\usepackage{tikz, tikz-cd}
\usepackage{pgf}
\usetikzlibrary{cd, babel}  % evita errores en tikz
\usetikzlibrary{graphs}
\usetikzlibrary{graphs.standard}
\usetikzlibrary{arrows,automata}

%% ADICIONALES Y GUÍAS, Y PARA MAYORES DETALLES VISITAR LA WEB: https://www.overleaf.com/learn/how-to/Creating_a_document_in_Overleaf

% GRAFICOS
\usepackage{graphicx}
% PONER CAPTIONS EN LAS IMAGENES
\usepackage[font=small,labelfont=bf]{caption}
% PONER SUBCAPTIONS EN LAS IMAGENES
\usepackage{subcaption}
% PONER LINKS
\usepackage{hyperref}
% PONER CÓDIGOS DE PROGRAMACION
\usepackage{multirow}
\usepackage{minted}
\usepackage{listings}
% PONER TÉRMINOS MATEMATICOS
\usepackage{amsmath}

%% ESTILO DE REFERENCIADO
\definecolor{LightGray}{gray}{0.9}
\usepackage[
backend=biber,
style=apa,      % Aquí, cambiar apa por el estilo de citado que requiera
sorting=ynt
]{biblatex}
\addbibresource{Referencias.bib}

%% ESQUEMA DE TABLA DE CONTENIDOS
\AtBeginSection[]
{
  \begin{frame}
    \frametitle{Contenidos}
    \tableofcontents[currentsection]
  \end{frame}
}
\AtBeginSubsection[]

%% Caratula / Información:
\title[Nombre Corto de la Presentación] {Nombre Completo de la Presentación}
\subtitle[Subtitulo]{Colocar el Subtitulo} % Opcional

% En caso de un autor único:
\author[N, Apellido1 (Año)]{Nombre Apellido1 Apellido2\texorpdfstring{\\}{, }correo-electronico@upc.edu.pe}
\institute[Facultad de Economía - UPC]{Facultad de Economía\\Universidad Peruana de Ciencias Aplicadas}

% En caso de más de un autor (Debe desactivar las lineas 71 y 72 colocando % al inicio, y eliminar los % de la linea 75 a 85):
%\author[ApellidoAutor1, ApellidoAutor2]{N.~ApellidoAutor1\inst{1} \and N.~ApellidoAutor2\inst{2}}
% \institute[FACUDECO] % Opcional, en funcion de cuantas afiliaciones existan.
%{
%  \inst{1,2}%
%  Facultad de Economía\\
%  Universidad Peruana de Ciencias Aplicadas
%  \and %Editar acorde con la afiliación de cada autor
%  \inst{2}%
%  School of Economics and Political Science\\
%  Very Famous University
% }
\date{XX de YYYYYY de 20ZZ} %(optional)

%%% LOS CONTENIDOS DEL DOCUMENTO SE REGISTRAN DESDE ESTA LÍNEA EN ADELANTE
\begin{document}

\begin{frame}
  \titlepage
\end{frame}

\section{Sección 1}

\begin{frame}
\frametitle{Titulo de la hoja 1}
\begin{itemize}
    \item Puede colocar texto en párrafos de manera continua. Aunque recuerde que se trata de una presentación, por lo que no se requiere de una gran cantidad de texto.
    \item En su defecto, se recomienda optar por el uso de viñetas.
    \item Asimismo, puede ir verificando si la hoja se encuentra recargada o si cuenta con suficiente información.
    \item Finalmente, no olvide marcar el inicio y final, como se puede ver en las líneas 95 y 103.
\end{itemize}
\end{frame}

\begin{frame}
\frametitle{Titulo de la hoja 2}
El titulo de la hoja puede variar, sin afectar el nombre de la sección en la que se encuentra.\medskip\par

\begin{enumerate}
    \item Nótese que la numeración se asigna de forma automática.
    \item Conforme se vaya familiarizando con el uso de Overleaf le será más sencillo incorporar otros tipos de elementos:
    \begin{itemize}
        \item Numeraciones.
        \item Viñetas simples.
        \item Subdivisiones, etc.
    \end{itemize}
    \item Para mayores detalles puede revisar la Guía para el caso de Presentaciones (conocido como Beamer): \hyperlink{}{https://es.overleaf.com/learn/latex/Beamer}
\end{enumerate}
\end{frame}
\section{Sección 2}

\begin{frame}
\frametitle{Titulo de la hoja 3}
Esta es una forma de colocar viñetas que vayan apareciendo en cada hoja o slide. No olvide colocar la numeración.
\begin{itemize}
 \item<1-> Texto visible en el 1er slide.
 \item<2-> Texto visible en el 2do slide.
 \item<3-> Texto visible en el 3er slide.
 \item<4-> Texto visible en el 4to slide.
\end{itemize}
\end{frame}
\section{Sección Final}

\begin{frame}
\frametitle{Titulo de la hoja 4}
\begin{itemize}
    \item Para mostrar tablas o ilustraciones, debe cargarlas como imágenes. Se sugiere usar la carpeta Imagenes (parte superior izquierda).
    \item Una vez cargada, puede usar la opción insertar en la opción de "..." de la barra superior.
    \item La carga de imágenes se verá como en el siguiente slide.
    \item Se sugiere guardar la imagen, con la footnote que detalle la fuente de origen de la imagen.
\end{itemize}
\end{frame}

\begin{frame}
\frametitle{Titulo de la hoja 4}
\begin{figure}
    \centering
    \caption{Titulo de la imagen o Tabla}    \includegraphics[width=0.6\linewidth]{Imagenes/Ejemplo1.png}
    \label{fig:Tabla1}
\end{figure}
\end{frame}

\begin{frame}
\frametitle{Titulo de la hoja 4}
\begin{itemize}
    \item Para colocar referencias, recuerde primero colocar los datos de la fuente usada en el archivo Referencias.bib.
    \item En el archivo encontrará ejemplos de referencias a Artículos Académicos y Documentos de Trabajo (\textit{article}), Libros (\textit{book}) y secciones de libros (\textit{inbook}). Pero existen otras más que puede indagar.
    \item Usada la información puede citarlos de las siguientes formas:
    \begin{itemize}
        \item Información basada en la revisión completa del autor (\cite{Einstein1905})
        \item Información basada en alguna parte del texto del autor (\cite[p.~150]{Knuth1973}).
        \item  O de la forma continua, como expresa Adrien y Dirac (\citeyear{Adrien&Dirac1981}).
    \end{itemize}
\end{itemize}
\end{frame}

\begin{frame}
  \titlepage
\end{frame}

\end{document}